\documentclass[10pt,a4paper]{amsart}
\usepackage[margin=19mm]{geometry}
\usepackage{amsmath,amssymb,mathtools}
\usepackage[scr=rsfs,cal=euler]{mathalfa}
\usepackage{xfrac}
\usepackage[unicode,hidelinks,bookmarks=false]{hyperref}
\newcommand{\ie}{i.e.\ }
\newcommand{\cf}{cf.\ }
\newcommand{\resp}{respectively\ }
\newcommand{\Subsection}{\subsection}
\providecommand{\nobreakdash}{\nobreak\hbox{-}}
\setlength{\emergencystretch}{2em}
\multlinegap0pt
\usepackage{bm}
\usepackage{bbm}
\usepackage{dsfont}
\usepackage{xspace}
\usepackage{cancel}
\usepackage{mathtools}
\usepackage{amsthm}
\makeatletter
\@ifundefined{theorem}{\newtheorem{theorem}{Theorem}}{}
\makeatother
\title[Bijectivity analysis of rational T-spline surfaces via Berns]{Bijectivity analysis of rational T-spline surfaces via Bernstein representations}
\author{Jia-Xuan Li, Ying-Ying Yu, Ya-Shu Liu, Xin Li, Ye Ji, Chun-Gang Zhu}
\date{Source: arXiv:2608.09838v1 (CC BY 4.0)}
\begin{document}
\maketitle

\noindent\emph{Source and licence.} Bijectivity analysis of rational T-spline surfaces via Bernstein representations. Version arXiv:2608.09838v1 (CC BY 4.0), distributed under the Creative Commons Attribution 4.0 International licence, as recorded in the arXiv licence metadata for this exact version and shown on the abstract page https://arxiv.org/abs/2608.09838v1. Full rights notice: licenses/arxiv-2608.09838-CC-BY-4.0.txt.\par\smallskip

\noindent\textit{Editorial scope.} Counted unit: the Theorem whose source label is \texttt{cor:regularity} in the article (source0.tex lines 506--522, source label \texttt{cor:regularity}), delivered together with the article's own complete original proof of it (source0.tex lines 524--583, one \texttt{proof} environment of 60 lines). No mathematics is written, repaired or re-derived: statement and proof are the article's own text line for line. Required context retained uncounted: thm:gram-bernstein (lines 420--432); eq:bernstein\_coefficients (lines 1363--1369). Every definition, display and label of those lines is retained unchanged. Omitted from this package and not redistributed: 1--419, 433--505, 523--523, 584--1362, 1370--1377. Nothing inside a retained excerpt is omitted or altered: every retained body is delivered exactly as the article's lines are written, so no omission receipt of any kind accompanies this package. Citations: the retained excerpts carry no citation command at all, so no bibliography entry of the article is transcribed and no reference is redistributed. Reference adaptation: none was needed and none was made; every reference inside the delivered unit resolves inside it, so no unresolved reference or citation is produced. Printed slips kept literal: no printed word, symbol, hypothesis or number of the counted statement or of its proof was altered. Layout and class: the article's own class is \texttt{jssc} with packages including lineno, hyperref, xcolor, algorithm2e, algorithmic, babel, bm, placeins, xfrac, epstopdf, cmap, amssymb, amsfonts, amsmath; the wrapper is the lane's pinned \texttt{amsart} editorial wrapper at 10pt on A4, which restates the theorem-like environments the retained text needs and adds the article's own macro definitions that the retained text needs (none), with extra wrapper packages (bm, bbm, dsfont, xspace, cancel, mathtools). The delivered numbering is the unit's own. Assets: no retained span contains a figure environment, a graphics-inclusion command, a PDF-page inclusion, a table or a TikZ picture; no image file is copied into this package, the package declares no asset dependency and no image file is redistributed with it. Licence: version arXiv:2608.09838v1 of this article is distributed under the Creative Commons Attribution 4.0 International licence, as recorded in the arXiv licence metadata for this exact version and shown on the abstract page https://arxiv.org/abs/2608.09838v1. A full rights notice is delivered at licenses/arxiv-2608.09838-CC-BY-4.0.txt. Characters: every retained excerpt is pure ASCII, so the delivered engine, XeLaTeX with its default Latin Modern text font, reports no missing character.
\begin{theorem}[Bernstein representation of the Gram determinant]
\label{thm:gram-bernstein}
Let $\bm{\mathcal{J}}^e(\xi,\eta)$ be the Jacobian matrix of a rational B\'ezier element $e$ of bi-degree $(p,q)$. Then the Gram determinant admits the Bernstein expansion
\begin{equation}
\det\!\left((\bm{\mathcal{J}}^e)^\top \bm{\mathcal{J}}^e \right)(\xi,\eta)
=
\frac{1}{(W^e(\xi,\eta))^{8}}
\sum_{r=0}^{8p-2}\sum_{s=0}^{8q-2}
D_{rs}^e \, B_r^{8p-2}(\xi) B_s^{8q-2}(\eta),
\label{eq:JeTJe-final}
\end{equation}
where the coefficients $D_{rs}^e\in\mathbb{R}$ depend only on the control points and weights associated with the element.
\end{theorem}

\begin{equation}
N^e(\xi,\eta)
=
\sum_{r=0}^{8p-2}\sum_{s=0}^{8q-2}
D_{rs}^e \, B_r^{8p-2}(\xi) B_s^{8q-2}(\eta),
\label{eq:bernstein_coefficients}
\end{equation}

\begin{theorem}[Sufficient condition for elementwise regularity]
\label{cor:regularity}
Assume that $W^e(\xi,\eta)>0$ on $[0,1]^2$. 
If all Bernstein coefficients in \eqref{eq:bernstein_coefficients} satisfy
\begin{equation}
D_{rs}^e \ge 0,
\qquad
\text{and at least one } D_{rs}^e > 0,
\end{equation}
then
\begin{equation}
\det\!\big((\bm{\mathcal{J}}^e)^\top\bm{\mathcal{J}}^e\big)(\xi,\eta)>0,
\qquad
(\xi,\eta)\in(0,1)^2.
\end{equation}
Consequently, $\bm{\mathcal{J}}^e(\xi,\eta)$ has full column rank on the interior of the element, and the parametrization $\bm{x}^e$ is regular on $(0,1)^2$.
\end{theorem}

\begin{proof}
By Theorem~\ref{thm:gram-bernstein}, the Gram determinant admits the representation
\begin{equation}
\det\!\big((\bm{\mathcal{J}}^e)^\top\bm{\mathcal{J}}^e\big)(\xi,\eta)
=
\frac{N^e(\xi,\eta)}{(W^e(\xi,\eta))^8},
\end{equation}
where
\begin{equation}
N^e(\xi,\eta)
=
\sum_{r=0}^{8p-2}\sum_{s=0}^{8q-2}
D_{rs}^e\, B_r^{8p-2}(\xi) B_s^{8q-2}(\eta).
\label{eq:numerator-Ne}
\end{equation}

We first establish the positivity of the numerator $N^e(\xi,\eta)$. 
For any $\xi\in(0,1)$ and any $r=0,\dots,8p-2$, the univariate Bernstein polynomial
\begin{equation}
B_r^{8p-2}(\xi)=\binom{8p-2}{r}\xi^r(1-\xi)^{8p-2-r}
\end{equation}
is strictly positive. Likewise, for any $\eta\in(0,1)$ and any $s=0,\dots,8q-2$,
\begin{equation}
B_s^{8q-2}(\eta)=\binom{8q-2}{s}\eta^s(1-\eta)^{8q-2-s}
\end{equation}
is strictly positive. Hence, each tensor-product basis function
\begin{equation}
B_r^{8p-2}(\xi)B_s^{8q-2}(\eta)
\end{equation}
is strictly positive on $(0,1)^2$.

Since $D_{rs}^e\ge 0$ for all $(r,s)$ and at least one coefficient is strictly positive, the expansion \eqref{eq:numerator-Ne} is a nonnegative linear combination of strictly positive functions with at least one strictly positive weight. Therefore,
\begin{equation}
N^e(\xi,\eta)>0,
\qquad
(\xi,\eta)\in(0,1)^2.
\end{equation}

Next, by assumption, $W^e(\xi,\eta)>0$ on $[0,1]^2$, and hence
\begin{equation}
(W^e(\xi,\eta))^8>0
\qquad\text{on }[0,1]^2.
\end{equation}
It follows that
\begin{equation}
\det\!\big((\bm{\mathcal{J}}^e)^\top\bm{\mathcal{J}}^e\big)(\xi,\eta)>0,
\qquad
(\xi,\eta)\in(0,1)^2.
\end{equation}

Finally, since $(\bm{\mathcal{J}}^e)^\top\bm{\mathcal{J}}^e$ is a $2\times 2$ Gram matrix, 
the positivity of its determinant implies that it is positive definite. 
Therefore, the column vectors of $\bm{\mathcal{J}}^e$ are linearly independent, and
\begin{equation}
\operatorname{rank}(\bm{\mathcal{J}}^e(\xi,\eta))=2,
\qquad
(\xi,\eta)\in(0,1)^2.
\end{equation}
This establishes the regularity of the parametrization on the interior of the element.
\end{proof}

\end{document}
